\section{Additional Details}
For all experiments, we freeze only the image tokenizer, DINOv2~\cite{oquab2023dinov2}, while fine-tuning all other components, including the frame-wise and global attention layers, camera and register tokens, and the point and camera prediction heads. We use a base learning rate of $2\!\times\!10^{-4}$ for all experiments, and apply a learning rate scaled by $1/10$ to the fine-tuned components inherited from the pretrained model. For evaluation on RE10K~\cite{zhou2018re10k} and ACID~\cite{liu2021acid}, we follow NoPoSplat and apply test-time camera pose optimization. The initial target camera parameters are obtained using the target-aligned projection formulation described in \cref{sec3.3:training_pipeline} of the main paper. For the ablation studies, we adopt a 5$\times$ reduced-iteration training schedule, where all iteration-related hyperparameters are reduced by a factor of 5. For the AnySplat results in \cref{fig:teaser} (a), we control the number of Gaussians by adjusting the voxel size. As the voxel size increases, we proportionally enlarge the scale of the output Gaussians so that their spatial coverage increases accordingly. 
